\documentclass[11pt]{article}
\usepackage[a4paper,margin=1in]{geometry}
\usepackage{booktabs}
\usepackage{array}
\usepackage{amsmath,amssymb}
\usepackage{xcolor}
\usepackage{fancyhdr}
\usepackage{enumitem}
\usepackage{longtable}
\usepackage[colorlinks=true,linkcolor=black,urlcolor=black,citecolor=black]{hyperref}

% AAAI-consolidated look: dark-blue section headings + running header rule
\definecolor{secblue}{RGB}{31,73,125}
\definecolor{boxgray}{RGB}{244,246,249}
\definecolor{noveltygreen}{RGB}{20,110,60}
\usepackage{sectsty}
\sectionfont{\color{secblue}\large}
\subsectionfont{\color{secblue}\normalsize}
\usepackage[most]{tcolorbox}

\pagestyle{fancy}
\fancyhf{}
\renewcommand{\headrulewidth}{0.4pt}
\lhead{\small Bruise Segmentation}
\rhead{\small Distillation Techniques \& Novelty Roadmap}
\cfoot{\thepage}

\newcolumntype{P}[1]{>{\raggedright\arraybackslash}p{#1}}
\setlength{\parindent}{0pt}
\setlength{\parskip}{5pt}

% Convenience macros
\newcommand{\Lsup}{\mathcal{L}_{\mathrm{sup}}}
\newcommand{\Lkd}{\mathcal{L}_{\mathrm{KD}}}
\newcommand{\Lfair}{\mathcal{L}_{\mathrm{fair}}}
\newcommand{\stopgrad}{\mathrm{sg}}
\newcommand{\novel}[1]{\textcolor{noveltygreen}{\textbf{#1}}}

\newtcolorbox{scopebox}{colback=boxgray,colframe=secblue,boxrule=0.6pt,
  arc=2pt,left=8pt,right=8pt,top=6pt,bottom=6pt}

\begin{document}

\begin{center}
  {\LARGE\bfseries Knowledge Distillation: Techniques and Novelty Roadmap}\\[4pt]
  {\large For the Fairness-Aware Bruise Segmentation Pipeline}\\[2pt]
  {\small SegFormer-B0 / B2 / B5 \, $\cdot$ \, YOLO26n-sem \, $\cdot$ \, paired ALS/WL \, $\cdot$ \, ITA skin-tone groups}
\end{center}

\vspace{2pt}

\begin{scopebox}
\textbf{What this document is.} A technical proposal that (i) synthesises the project's own
nine-technique distillation catalogue with the 34-paper fairness-aware KD literature review,
and (ii) proposes \novel{novel mix-and-match techniques} not previously combined for bruise
segmentation. \textbf{No new numbers are claimed} --- every technique here is a \emph{candidate}
for experiments. Existing results are quoted only as anchors. All comparisons are assumed to run
through the project's established protocol (paired subject-level cluster bootstrap, 28 subjects,
$B{=}4000$, median-Dice MDE $\approx 0.04$), and every claim is framed on
\textbf{complete-miss rate and worst-group performance}, not mean Dice, because the task is
\textbf{label-limited, not capacity-limited} (seven architectures cluster at 0.74--0.79;
human--human agreement $\approx 0.639$).
\end{scopebox}

\section{Executive Summary}

The single most valuable novelty available to this project is already \emph{half-observed}:
\textbf{the distilled student beats its teacher}. On the 185-image test set, SegFormer-B0
\emph{distilled} reaches mean Dice \textbf{0.786} against the SegFormer-B2 teacher's
\textbf{0.773}, at $7.4\times$ fewer parameters (3.71\,M vs 27.35\,M). A student that
outperforms a $7\times$ larger teacher is the strongest possible framing for a distillation
paper --- but at present it is a single-condition observation inside the $\sim$0.04 MDE. The
techniques below are organised around \emph{turning that anecdote into a robust, explained,
fairness-audited result}.

Three groups of ideas follow:
\begin{enumerate}[leftmargin=1.4em,itemsep=1pt]
  \item \textbf{Catalogue (\S4)} --- nine architecture-matched techniques from the project's
  own proposal doc, each now cross-referenced to the fairness-KD literature that validates it.
  \item \textbf{Novel combinations (\S5)} --- \novel{nine mix-and-match techniques} whose
  specific combination does not appear in the bruise-segmentation or fairness-KD literature.
  These are the paper-worthy contributions.
  \item \textbf{Student-beats-teacher playbook (\S6)} and \textbf{beyond distillation (\S7)}.
\end{enumerate}

\section{Project Assets and What Already Exists}

\begin{center}
\small
\begin{tabular}{P{3.1cm} P{4.0cm} P{4.4cm} P{2.0cm}}
\toprule
Asset & Detail & Role in distillation & Params (M) \\
\midrule
SegFormer-B5 & largest 2021 MiT family member & \emph{candidate} big teacher (gated) & $\sim$82--85 \\
SegFormer-B2 & current teacher & response-KD teacher (WL, ALS) & 27.35 \\
SegFormer-B0 & current student & distilled + direct & 3.71 \\
YOLO26n-sem & fastest model, high miss & pseudo-label student; DML candidate & 1.63 \\
Paired ALS/WL & same bruise, two modalities & cross-modal teacher (ALS$\to$WL) & --- \\
ITA groups (5) & per-image skin-tone index & fairness reweighting / gating & --- \\
\bottomrule
\end{tabular}
\end{center}

\textbf{Already implemented (baselines, not to be re-proposed).} Online response-based KD
(WL$\to$WL, median Dice 0.828); offline pseudo-label KD (YOLO); cross-modal KD (ALS$\to$WL,
median 0.811); fairness-reweighted KD (Approach A/B, the latter adds an MMD feature term).
\emph{Every one of these matches only the final output probability.} The gap --- intermediate
features, spatial structure, multi-source fusion, and miss-rate-targeted transfer --- is what
the new techniques address.

\section{Technique Catalogue, Cross-Referenced to the Literature}

Each row is a technique from the project's internal proposal, now paired with the
fairness-KD paper(s) that independently support it. ``Cost'' is implementation + compute
relative to the existing response-KD path.

\begin{center}
\small
\begin{longtable}{P{2.7cm} P{4.6cm} P{4.4cm} P{1.5cm}}
\toprule
Technique & What it transfers & Literature anchor & Cost \\
\midrule
\endhead
Feature / hint (FitNets) & intermediate encoder feature maps via $1{\times}1$ adapters & ClearKD (Huang 2024) refines teacher features for medical KD & small \\
Attention transfer & single-channel spatial attention per stage; no adapter & spatial-focus transfer; complements hint & near-zero \\
Structured / channel-wise & per-channel spatial distributions (KL) + pixel-affinity (MSE) & scale-invariant relational KD & small \\
Boundary-aware KD & KD loss reweighted by distance-to-edge & segmentation-fairness gap (Ben\v{c}evi\'c 2023) & near-zero \\
Multi-teacher ensemble & WL-teacher $+$ ALS-teacher fused soft target & AGRE-KD (Kenfack 2024); ensemble-KD (Radwan 2024) & moderate \\
Deep mutual learning & SegFormer-B0 $\leftrightarrow$ YOLO, live, no frozen teacher & DSP-KD progressive transfer (Zeng 2024) & moderate \\
Self-distillation & network's own deep layers teach shallow (raises ceiling) & last-layer transplant debiasing (Lee 2023) & small \\
Adversarial fairness & GRL makes skin tone unrecoverable from features & adversarial debiasing (2023); FairDisCo (2022) & moderate \\
Compression-aware & INT8 / pruned student, KD-recovered & FairPrune (Wu 2022); channel pruning (Kong 2024) & small \\
\bottomrule
\end{longtable}
\end{center}

\noindent\textbf{Generic response-KD loss (the shared base).}
\[
\mathcal{L} \;=\; \alpha\,\Lsup(z_s, y) \;+\; (1-\alpha)\,\mathrm{BCE}\!\big(z_s,\ \sigma(z_t/T)\big),
\]
temperature-scaled, calibrated soft targets (Menon-style, as in
\texttt{pipeline/losses.py::DistillSegLoss}). Every technique below is an
\emph{additive term} or a \emph{reweighting} of this base.

\section{Novel Mix-and-Match Techniques}

The literature review's central finding is that fairness-KD methods
(FairDi, AGRE-KD, two-biased-teachers, FairPrune, FairDisCo, BMFT, causal modeling)
were developed almost entirely for \textbf{classification}. Segmentation fairness is called out
as ``underexplored, existing methods ineffective'' (Ben\v{c}evi\'c 2023; the review's
\emph{Fairness in Segmentation} gap, priority Medium). \novel{Porting these methods to
pixel-level bruise segmentation, and fusing them with this project's unique
ALS/WL pairing and complete-miss metric, is where the novelty lives.} Mix-and-match counts as
novelty when the specific combination has not appeared before --- each idea below is annotated
with why it is new.

\begin{center}
\small
\begin{longtable}{@{}p{0.4cm} P{5.3cm} P{5.5cm} c@{}}
\toprule
\# & Novel technique & Why it is new (not in either source) & Priority \\
\midrule
\endhead
N1 & \novel{Complete-miss-aware distillation} & KD optimises Dice/logits everywhere; nobody weights the loss by \emph{teacher-confident $\wedge$ student-missed} regions to kill the project's headline failure (YOLO 7--9\% miss) & High \\
N2 & \novel{Heterogeneous cross-family $+$ cross-modal ensemble} & B5(WL) $+$ B2(ALS) $\to$ B0: teachers differ in \emph{both} architecture family \emph{and} imaging modality; ensemble-KD papers use same-modality same-family teachers & High \\
N3 & \novel{Boundary $\times$ fairness joint reweighting} & one per-pixel weight $= w_{\text{edge}}\!\cdot w_{\text{group}}$; boundary-KD and group-KD both exist, their product does not & High \\
N4 & \novel{ALS-oracle uncertainty-gated KD} & use ALS soft targets \emph{only} where the WL teacher is uncertain --- ALS reveals bruises invisible on WL, disproportionately on darker skin & High \\
N5 & \novel{Group-robust ensemble weighting for segmentation} & AGRE-KD's adaptive per-group teacher weights, ported from classification to pixel-level, gated by the surviving light-skin recall deficit & Medium \\
N6 & \novel{Ceiling-raising self-distillation $\to$ response KD} & self-distil the B5/B2 teacher first (raise the ceiling), \emph{then} distil to B0 --- chaining the only technique that improves the teacher with the transfer step & Medium \\
N7 & \novel{Disagreement-mined mutual learning} & SegFormer-B0 $\leftrightarrow$ YOLO co-training where inter-model disagreement \emph{is} the hard-example signal (fuses DML with hard-example mining) & Medium \\
N8 & \novel{Pixel-level adversarial fairness (GRL)} & FairDisCo/adversarial debiasing are image-level classifiers; a per-pixel GRL on decode features conditioned on ITA is untried for bruise seg & Medium \\
N9 & \novel{Fairness-preserving compression} & INT8/prune B0 with KD recovery while \emph{monitoring per-group miss-rate} --- FairPrune is classification; segmentation edge-deploy fairness is open & Low \\
\bottomrule
\end{longtable}
\end{center}

\subsection{N1 --- Complete-Miss-Aware Distillation \novel{(top pick)}}
The project's own reporting leads with complete-miss rate, not Dice --- YOLO misses 7--9\% of
bruises entirely, and misses concentrate on hard cases (small, low-contrast, and plausibly
specific skin tones). Standard KD spends equal loss budget on every pixel and so barely moves
the miss rate. Define a per-pixel weight that fires exactly where the teacher is confident there
is a bruise but the student predicts background:
\[
w_{\text{miss}}(x,y) \;=\; 1 \;+\; \gamma\,\underbrace{p_t(x,y)}_{\text{teacher conf.}}\cdot
\underbrace{\big(1-\sigma(z_s(x,y))\big)}_{\text{student misses}},\qquad
\mathcal{L}^{\text{miss}}_{\text{KD}} = \operatorname*{mean}_{x,y}\big[w_{\text{miss}}\cdot\mathrm{BCE}(z_s,p_t)\big].
\]
\textbf{Why novel:} KD losses in both sources optimise agreement uniformly; none target the
\emph{complete-miss} event, which is this project's primary metric. \textbf{Hook:} one-line
reweight of \texttt{DistillSegLoss}, reusing the teacher forward already in the loop. \textbf{Risk:}
low --- reduces to standard KD at $\gamma{=}0$.

\subsection{N2 --- Heterogeneous Cross-Family + Cross-Modal Ensemble \novel{(top pick)}}
Train one WL-only B0 student against a single soft target fused from \emph{two structurally
different} teachers: SegFormer-B5 on the WL image and SegFormer-B2 on the paired ALS image.
\[
p_{\text{ens}}(x,y) \;=\; g(x,y)\,p^{\,\text{B5,WL}}(x,y) \;+\; \big(1-g(x,y)\big)\,p^{\,\text{B2,ALS}\to\text{WL}}(x,y).
\]
The gate $g$ can be fixed, confidence-based, or conditioned on ITA group. \textbf{Why novel:}
ensemble-KD papers (Radwan, Kenfack) combine same-modality, same-family teachers; here the two
teachers disagree by \emph{design} --- different capacity family \emph{and} different physics
(ALS fluorescence vs white light). \textbf{Gate:} only worthwhile if the Phase-1 disagreement
analysis shows B5/B2 make complementary errors; otherwise single-teacher B5 KD suffices.

\subsection{N3 --- Boundary $\times$ Fairness Joint Reweighting}
Both a boundary weight (from the distance transform already in
\texttt{metrics\_extended.py}) and a group weight (from the ITA label already used by the
fairness losses) can multiply the same pixel-wise KD term:
\[
w(x,y)=\underbrace{\exp\!\big(-d(x,y)^2/2\sigma^2\big)}_{\text{edge focus}}\cdot
\underbrace{\lambda_{g(x,y)}}_{\text{ITA-group weight}},\qquad
\mathcal{L} = \operatorname*{mean}\big[w\cdot\mathrm{BCE}(z_s,p_t)\big].
\]
\textbf{Why novel:} each factor exists in isolation; their product --- ``sharpen boundaries
\emph{most} for the worst-served skin-tone group'' --- does not appear in either source and
directly serves the project's surface-Dice/HD95 \emph{and} fairness goals at once.

\subsection{N4 --- ALS-Oracle Uncertainty-Gated KD}
ALS (alternate light source) fluorescence reveals bruising that is invisible or faint under
white light, and the clinical motivation for the whole project is that this gap is largest on
\emph{darker} skin. Instead of a fixed ALS$\to$WL distillation, use the ALS teacher as a
\emph{selective oracle}: trust ALS soft targets only where the WL teacher is uncertain.
\[
p^{*}(x,y)=\big(1-u(x,y)\big)\,p^{\text{WL}} + u(x,y)\,p^{\text{ALS}},\quad
u(x,y)=\mathbf{1}\!\big[\,|p^{\text{WL}}-\tfrac12| < \tau\,\big],
\]
with $u$ the WL-teacher uncertainty gate. \textbf{Why novel:} cross-modal KD exists in the
project, but \emph{uncertainty-gated, fairness-motivated} modality fusion --- ALS as a targeted
corrector for exactly the ambiguous, disproportionately dark-skin regions --- is new and
clinically grounded.

\subsection{N5 --- Group-Robust Ensemble Weighting for Segmentation}
Port AGRE-KD (Kenfack 2024, adaptive group-robust ensemble KD, classification) to
pixel-level: set the ensemble gate weights per ITA group from \emph{validation worst-group
recall}, so the ensemble leans on whichever teacher best serves the currently-worst-served
group. \textbf{Why novel:} AGRE-KD has never been applied to segmentation, and the group signal
is anchored to the one fairness quantity that has survived clustering here (the light-skin
recall deficit).

\subsection{N6 --- Ceiling-Raising Self-Distillation, then Transfer}
Every other technique assumes the teacher is fixed. Self-distillation (auxiliary heads on early
encoder stages, each matching GT \emph{and} the network's own deepest stop-grad output) is the
one technique that improves the \emph{teacher itself}. Apply it to B5/B2 first, then run N1--N4
from the improved teacher. \textbf{Why novel here:} chaining ceiling-raising self-distillation
into a downstream fairness-aware transfer is a two-stage recipe neither source proposes; it also
gives the cleanest ``student beats an even-stronger teacher'' story (\S6).

\subsection{N7 --- Disagreement-Mined Mutual Learning}
Deep mutual learning trains SegFormer-B0 and YOLO26n-sem together, each the other's live soft
target (no frozen teacher). Novelty: use the \emph{per-pixel inter-model disagreement} as an
automatic hard-example weight --- the two very different architectures disagree exactly on the
ambiguous cases that drive the annotation ceiling and the miss rate. \textbf{Caveat:} needs a
supervised warmup to avoid joint collapse; a genuinely new dual-optimizer training script.

\subsection{N8 --- Pixel-Level Adversarial Fairness}
A gradient-reversal discriminator tries to predict ITA group from the student's decode features;
the backbone is pushed to make skin tone unrecoverable. FairDisCo and adversarial debiasing do
this for image-level classifiers; \novel{a per-pixel GRL for a segmentation decode head} is
untried for bruise work. Frame as ``Approach C,'' run \emph{alongside} the existing MMD Approach
B for comparison, since adversarial training is less stable than moment matching.

\subsection{N9 --- Fairness-Preserving Compression}
Prune/INT8-quantise the distilled B0 with a short KD recovery fine-tune (teacher = full-precision
B0), \emph{gated on per-ITA-group miss-rate not dropping}. FairPrune (Wu 2022) and channel
pruning (Kong 2024) show pruning can \emph{improve} fairness in classification; whether that
holds for segmentation edge deployment is open. Only worthwhile once a sub-10\,ms on-device
target is real (B0 already hits 86 FPS on GPU).

\section{The Student-Beats-Teacher Playbook}

This is the headline the user asked to prioritise, and the data already leans in its favour
(B0-distilled mean Dice $0.786 > 0.773$ B2). To make it a \emph{robust, defensible} claim rather
than a within-MDE coincidence:

\begin{enumerate}[leftmargin=1.4em,itemsep=2pt]
  \item \textbf{Establish it with statistics, not a point estimate.} Report the paired
  student$-$teacher difference with its 95\% cluster-bootstrap CI on \emph{mean Dice,
  complete-miss rate, and worst-ITA-group recall} --- not median Dice, where the gap is inside
  the MDE. ``Beats on mean Dice and matches on median at $7.4\times$ smaller'' is already true
  and publishable.
  \item \textbf{Explain \emph{why} a student can beat its teacher.} The mechanism here is
  \emph{regularisation}: soft teacher targets smooth the label noise that caps every model at the
  annotation ceiling. A student trained on calibrated soft targets can generalise past a teacher
  trained on hard, noisy masks. This reframes ``student beats teacher'' from a curiosity into
  \emph{evidence that the task is label-limited} --- the project's central thesis.
  \item \textbf{Strengthen the effect deliberately} with the techniques above, in expected order
  of payoff: N1 (miss-aware) and N4 (ALS-oracle) attack the exact failures the teacher makes;
  N2 (heterogeneous ensemble) gives the student a target better than \emph{any single} teacher,
  which is the cleanest route to a student that beats every individual teacher; N6 raises the
  teacher ceiling so the win is against a \emph{stronger} baseline.
  \item \textbf{Audit fairness of the win.} A student that beats the teacher on mean Dice but
  worsens the worst-group deficit is not a real win. Report the per-group breakdown every time.
\end{enumerate}

\section{Beyond Distillation --- What Else We Can Do}

The literature review flags gaps that distillation alone will not close; these are
complementary contributions the project is well-positioned for.

\begin{center}
\small
\begin{tabular}{P{3.6cm} P{7.2cm} P{2.6cm}}
\toprule
Direction & What and why & Fit \\
\midrule
Skin-tone scale robustness & Report fairness under Monk / chromatic-clustering \emph{as well as} ITA; Desai (2025) + Kalb (2023) show ITA is high-variance and granular scales expose disparities the project may be masking & High (eval-only, cheap) \\
Minority-tone augmentation & Diffusion / image-to-image synthesis (FairSkin; Wang majority$\to$minority) to enlarge under-represented ITA groups --- directly attacks the annotation ceiling & Medium \\
Active annotation on misses & Semi-supervised / active learning targeting complete-miss images; the real bottleneck is labels, not capacity & High (addresses root cause) \\
Calibration-aware eval & Temperature-fit every model and report ECE per group; soft-target KD only helps if teachers are calibrated & Medium \\
Multi-attribute fairness & Move beyond single-attribute ITA toward intersection with bruise size/age --- the review's ``lack of multi-attribute fairness'' gap & Low \\
\bottomrule
\end{tabular}
\end{center}

\section{Prioritised Roadmap}

Ordered by (payoff) $\div$ (implementation cost $\times$ risk), reusing existing code wherever
possible and respecting the ``ablate one component at a time'' rule.

\begin{center}
\small
\begin{tabular}{@{}c P{5.0cm} P{4.6cm} P{2.6cm}@{}}
\toprule
Order & Do this & Because & New infra? \\
\midrule
1 & N1 complete-miss-aware KD & one-line reweight; targets the headline metric; cannot hurt at $\gamma{=}0$ & no \\
2 & N3 boundary $\times$ fairness reweight & reuses distance-transform + ITA weights already present & no \\
3 & Phase-1 gate: B5 vs B2 disagreement & decides whether N2/N5 ensembles are worth it & no \\
4 & N2 heterogeneous ensemble (if gate passes) & cleanest ``beats every single teacher'' result & merges 2 loops \\
5 & N4 ALS-oracle gated KD & clinically grounded fairness lever; ALS teacher exists & moderate \\
6 & N6 self-distil teacher, then re-distil & raises ceiling; strengthens \S6 story & new head \\
7 & N5 group-robust ensemble weighting & fairness refinement of N2 & small \\
8 & N8 adversarial fairness (Approach C) & compare vs MMD Approach B & moderate \\
9 & N7 DML, N9 compression & genuinely new training infra; deployment-gated & yes \\
\bottomrule
\end{tabular}
\end{center}

\begin{scopebox}
\textbf{Headline experiment if resources force one choice:}
\novel{B5$\to$B0 miss- and fairness-aware distillation} ---
$\mathcal{L} = \alpha\Lsup + (1-\alpha)\,\mathcal{L}^{\text{miss}}_{\text{KD}}
+ \lambda_{\text{edge$\times$grp}}\mathcal{L}^{\text{N3}}$, extended to the heterogeneous
B5(WL)$+$B2(ALS) ensemble (N2) only if the Phase-1 gate shows complementary errors.
Report on complete-miss rate and worst-ITA-group recall with cluster CIs; never run the whole
stack at once.
\end{scopebox}

\section*{\textcolor{secblue}{Key References (Fairness-Aware KD Literature)}}
\small
\begin{itemize}[leftmargin=1.2em,itemsep=1pt]
\item Masroor et al.\ (2024). \emph{Fair Distillation: teaching fairness from biased teachers in medical imaging.}
\item Kenfack et al.\ (2024). \emph{Adaptive Group-Robust Ensemble Knowledge Distillation (AGRE-KD).}
\item Pundhir et al.\ (2024). \emph{Biasing \& debiasing (two-biased-teachers) for equitable skin analysis.}
\item Zeng et al.\ (2024). \emph{DSP-KD: dual-stage progressive KD for skin disease classification.}
\item Yu et al.\ (2024). \emph{Variational AdaBoost KD for skin lesion classification.}
\item Huang et al.\ (2024). \emph{ClearKD: clear KD for medical image classification.}
\item Lee \& Lee (2023). \emph{Debiased distillation by transplanting the last layer.}
\item Wu et al.\ (2022). \emph{FairPrune: fairness through pruning.} \quad Kong et al.\ (2024). \emph{Channel pruning for fairness.}
\item Xue et al.\ (2024). \emph{BMFT: bias-based weight masking fine-tuning.}
\item \emph{FairDisCo} (2022); Adversarial debiasing for fair skin lesion classification (2023).
\item Ben\v{c}evi\'c et al.\ (2023). \emph{Understanding skin color bias in deep learning-based skin lesion segmentation.} (segmentation gap)
\item Desai et al.\ (2025). \emph{The effect of skin tone classification on bias in bruise detection.} (YOLOv5; ITA vs granular scales)
\item Kalb et al.\ (2023). \emph{Revisiting skin tone fairness} (ITA variability). \quad Zhang et al.\ (2024). \emph{FairSkin: fair diffusion.}
\end{itemize}

\vfill
\begin{center}\footnotesize\itshape
Proposal document --- no results claimed. Every technique is a candidate for future experimentation,
to be validated on complete-miss rate and worst-group performance under the project's paired
cluster-bootstrap protocol.
\end{center}

\end{document}
